 \vspace{-4mm}
\begin{figure}[!ht]
    \centering
    \includegraphics[width=\textwidth]{figs/teaser.pdf}
   \vspace{-6mm}
    \caption{\textbf{(a)} Comparisons of perception-based and latent world model-based approaches on nuScenes and NavSim benchmarks. \textbf{(b)} Planning visualization on the front view and bird's-eye-view (BEV) space. Our DynFlowDrive achieves comparable performance.}
    \vspace{-8mm}
    \label{fig:teaser} %
\end{figure}

\begin{abstract}

Recently, world models have been incorporated into the autonomous driving systems to improve the planning reliability. Existing approaches typically predict future states through appearance generation or deterministic regression, which limits their ability to capture trajectory-conditioned scene evolution and leads to unreliable action planning.
To address this, we propose \textbf{DynFlowDrive}, a latent world model that leverages flow-based dynamics to
model the transition of world states under different driving actions. By adopting the rectified-flow formulation, the model learns a velocity field that describes how the scene state changes under different driving actions, enabling progressive prediction of future latent states.
Building upon this, we further introduce a stability-aware multi-mode trajectory selection strategy that evaluates candidate trajectories according to the stability of the induced scene transitions.
Extensive experiments on the nuScenes and NavSim benchmarks demonstrate consistent improvements across diverse driving frameworks without introducing additional inference overhead. Source code will be abaliable at \href{https://github.com/xiaolul2/DynFlowDrive}{https://github.com/xiaolul2/DynFlowDrive}.

\vspace{-2mm}
\keywords{Autonomous Driving \and Latent World Model \and Trajectory Planning}


\end{abstract}
